%% Supplementary material for "Recognising the phone, not the throat" (Computer Methods and Programs in Biomedicine).
%% Compile: tectonic supplementary.tex. Tables in supp_tables/ are generated by `python -m src.supp_tables`.
\documentclass[11pt]{article}
\usepackage[a4paper,margin=2.2cm]{geometry}
\usepackage{amsmath,amssymb}
\usepackage{graphicx}
\usepackage{booktabs}
\usepackage{array}
\usepackage{longtable}
\usepackage[hidelinks]{hyperref}
\usepackage{xurl}
\usepackage[numbers,sort&compress]{natbib}
\setlength{\emergencystretch}{1.5em}
\renewcommand{\thefigure}{S\arabic{figure}}
\renewcommand{\thetable}{S\arabic{table}}
\renewcommand{\thesection}{S\arabic{section}}
\renewcommand{\theequation}{S\arabic{equation}}
\newcommand{\figfile}[2]{\IfFileExists{#1}{\includegraphics[width=#2]{#1}}{\fbox{\parbox{0.8\linewidth}{\centering\vspace{2cm}Missing file: \texttt{\detokenize{#1}}\vspace{2cm}}}}}

\title{Supplementary material\\[4pt]\large Recognising the phone, not the throat: deconfounded evaluation and shortcut mitigation for smartphone pharyngitis classification}
\author{Sandesh Bhatta\\\small Department of Computer Science, University at Albany, State University of New York}
\date{}

\begin{document}
\maketitle
\tableofcontents

%% -------------------------------------------------------------------------------------------
\section{Differences between the two phones}

Table~\ref{tab:s-symptoms} compares the symptom records of patients photographed with each phone (chi-square tests, Bonferroni correction over 20 symptoms). Samsung patients reported more symptoms (mean 2.7 vs 1.7, Mann--Whitney $p < 10^{-19}$), mostly systemic ones, yet had a lower bacterial prevalence. Labelling also differed (main text, Table~1): Samsung patients were rated by more physicians (median 5 vs 4; only Samsung patients had 8 or 9 raters), physicians agreed more often (mean agreement 0.66 vs 0.47, $p < 10^{-11}$), and ties were rarer (6.3\% vs 23.1\%, $p < 10^{-9}$).

\begin{table}[!htbp]
\centering
\small
\caption{Share of patients with each symptom, by phone (\%), ordered by $p$ value.}
\label{tab:s-symptoms}
\begin{tabular}{@{}lcccc@{}}
\toprule
Symptom & Samsung ($n = 382$) & Xiaomi ($n = 290$) & $p$ & $p$ (Bonferroni) \\
\midrule
\csname @@input\endcsname supp_tables/symptoms.tex 
\bottomrule
\end{tabular}
\end{table}

%% -------------------------------------------------------------------------------------------
\section{Validation of the deconfounded AUC: full results}

Table~\ref{tab:s-sim} gives the simulation results summarised in main-text Fig.~3A (500 replicates per setting, 700 patients, two sites, overall prevalence 0.28). The \textit{target} is the AUC of the same score in a large unconfounded population from the same sites (prevalence gap 0). The deconfounded AUC was unbiased for this target in every setting (largest absolute bias 0.001). The adjusted AUC was unbiased for the marker-only and site-only scores but overestimated the target of the learned model by up to 0.024, because within-site comparisons ignore the site offset that the learned model adds to its scores. Table~\ref{tab:s-sweep} gives the resampling results on PGUPharyngitis summarised in main-text Fig.~3B (20 resamples per gap, 200 patients per phone, overall prevalence 0.25).

\begin{table}[!htbp]
\centering
\footnotesize
\caption{Simulation. Mean metric (mean bias against the target) over 500 replicates per setting; last column: SD of the deconfounded / adjusted AUC across replicates.}
\label{tab:s-sim}
\setlength{\tabcolsep}{3pt}
\begin{tabular}{@{}lcccccc@{}}
\toprule
Score & Gap & Target & Pooled & Deconfounded & Adjusted & SD \\
\midrule
\csname @@input\endcsname supp_tables/simulation.tex 
\bottomrule
\end{tabular}
\end{table}

\begin{table}[!htbp]
\centering
\footnotesize
\caption{PGUPharyngitis resampled to a given prevalence gap between phones. Mean (SD) over 20 resamples, each evaluated with stratified 5-fold cross-validation. At gap 0 the phone carries no label information.}
\label{tab:s-sweep}
\begin{tabular}{@{}lcccc@{}}
\toprule
Model & Gap & Pooled AUC & Deconfounded AUC & Adjusted AUC \\
\midrule
\csname @@input\endcsname supp_tables/sweep.tex 
\bottomrule
\end{tabular}
\end{table}

%% -------------------------------------------------------------------------------------------
\section{Mitigation: adjusted AUC and cross-phone transfer}

Table~\ref{tab:s-mit} adds the covariate-adjusted AUC \citep{janes2008adjusting} to main-text Table~5. No mitigation method changed the adjusted AUC significantly (corrected $t$-test, all $p > 0.10$). Table~\ref{tab:s-cross} gives cross-phone transfer with colour-constancy features and per-phone standardisation (each phone standardised with its own statistics, which needs no labels from the test phone).

\begin{table}[!htbp]
\centering
\footnotesize
\caption{Pooled, deconfounded and adjusted AUC of the mitigation methods (frozen embeddings, 5\,$\times$\,5 cross-validation, mean over repeats).}
\label{tab:s-mit}
\setlength{\tabcolsep}{3.5pt}
\begin{tabular}{@{}>{\raggedright\arraybackslash}p{4.2cm}cccccc@{}}
\toprule
& \multicolumn{3}{c}{DINOv2} & \multicolumn{3}{c}{ConvNeXt-Tiny} \\
\cmidrule(lr){2-4}\cmidrule(l){5-7}
Method & Pooled & Deconf. & Adjusted & Pooled & Deconf. & Adjusted \\
\midrule
\csname @@input\endcsname supp_tables/mitigation_adjusted.tex 
\bottomrule
\end{tabular}
\end{table}

\begin{table}[!htbp]
\centering
\footnotesize
\caption{Cross-phone transfer (AUC [95\% bootstrap CI]): trained on all patients of one phone, tested on the other.}
\label{tab:s-cross}
\begin{tabular}{@{}llcc@{}}
\toprule
Features & Correction & Samsung $\rightarrow$ Xiaomi & Xiaomi $\rightarrow$ Samsung \\
\midrule
\csname @@input\endcsname supp_tables/cross.tex 
\bottomrule
\end{tabular}
\end{table}

%% -------------------------------------------------------------------------------------------
\section{Gated cross-modal fusion network}

The image embedding passes through layer normalisation, dropout and a linear layer to a 128-dimensional vector $h_{\text{img}}$; the 22 tabular inputs pass through a 22--64--128 perceptron to $h_{\text{sym}}$. A gate computes a scalar weight per patient,
\begin{equation}
g = \sigma\!\left(W [h_{\text{img}}; h_{\text{sym}}] + b\right), \qquad z = g\, h_{\text{img}} + (1-g)\, h_{\text{sym}},
\end{equation}
and a dropout--linear head maps $z$ to a logit. Auxiliary image-only and symptom-only heads add their losses with weight $\lambda$, and modality dropout (probability 0.2 per branch, never both) randomly removes one branch during training. Training used the soft label (share of physicians voting bacterial) as target in class-weighted binary cross-entropy, AdamW with weight decay $10^{-4}$, batch 32 and a cosine schedule over 60 epochs. Each epoch drew one of the ten cached augmentation embeddings per patient. The inner cross-validation selected the learning rate $\{10^{-3}, 3 \times 10^{-4}\}$, dropout $\{0.3, 0.5\}$, $\lambda \in \{0, 0.3\}$ and the number of epochs (maximum of the smoothed mean inner AUC curve); the model was then retrained on the full training part. Ablations replaced the gate by concatenation or by FiLM \citep{perez2018film}, used hard labels, removed the auxiliary heads, removed modality dropout, or removed fever and cough from the inputs.

\paragraph{Ablations} No component made a measurable difference (Fig.~\ref{fig:s-ablation}). With ConvNeXt features, the full gated model reached 0.684; replacing the gate by concatenation gave 0.686 ($p = 0.60$), FiLM 0.681, removing the auxiliary heads 0.682 and removing modality dropout 0.682. Training on hard instead of soft labels lowered AUC to 0.672 (difference 0.012, $p = 0.19$). Removing fever and cough did not hurt (0.690). A shuffled-image control, in which image embeddings were permuted across patients, fell to 0.529, below symptoms alone, so the pipeline does not leak labels.

\paragraph{Gate} The gate weighted the image at $g \approx 0.48$ on average (DINOv2) and slightly less when physicians disagreed (median 0.47 vs 0.50, Mann--Whitney $p < 0.001$; Fig.~\ref{fig:s-gate}).

\begin{figure}[!htbp]
\centering
\figfile{figures/fig4_ablation.pdf}{0.7\textwidth}
\caption{Ablation of the gated fusion network (ConvNeXt features): mean fold AUC with 95\% confidence interval. The shuffled-image control falls to chance.}
\label{fig:s-ablation}
\end{figure}

\begin{figure}[!htbp]
\centering
\figfile{figures/figS3_gate.pdf}{\textwidth}
\caption{Gate values of the gated fusion model (DINOv2): distribution for unanimous and split votes (left) and median gate value by number of positive symptoms (right). $g = 1$ means the model relies only on the image.}
\label{fig:s-gate}
\end{figure}

%% -------------------------------------------------------------------------------------------
\section{Planned fusion hypotheses under three tie rules}

Gated fusion (DINOv2) beat the best symptom model by 0.095 AUC (Holm-adjusted $p = 0.012$; significant in 5 of 5 DeLong repeats) but did not beat the best image-only model ($-0.020$, $p = 1.0$) or the best standard fusion ($-0.008$, $p = 1.0$). Under the two alternative tie rules the advantage over symptoms was no longer significant (Table~\ref{tab:s-hyp}).

\begin{table}[!htbp]
\centering
\footnotesize
\caption{Planned comparisons for gated fusion (DINOv2). $\Delta$AUC is the mean fold-level difference; $p$ values are one-sided, Nadeau--Bengio corrected and Holm-adjusted.}
\label{tab:s-hyp}
\setlength{\tabcolsep}{4pt}
\begin{tabular}{@{}>{\raggedright\arraybackslash}p{4.2cm}ccc@{}}
\toprule
Comparison & Ties bacterial & Ties non-bacterial & Ties dropped \\
\midrule
H1: vs best symptom-only & $+0.095$ ($p = 0.012$) & $+0.044$ ($p = 0.38$) & $+0.051$ ($p = 0.22$) \\
H2: vs best image-only & $-0.020$ ($p = 1.0$) & $-0.009$ ($p = 1.0$) & $-0.002$ ($p = 1.0$) \\
H3: vs best standard fusion & $-0.008$ ($p = 1.0$) & $-0.018$ ($p = 1.0$) & $-0.019$ ($p = 1.0$) \\
\addlinespace
Image (DINOv2): AUC / within phone & 0.700 / 0.586 & 0.671 / 0.609 & 0.696 / 0.605 \\
Phone only: AUC & 0.673 & 0.618 & 0.650 \\
Best symptom-only: AUC & 0.596 & 0.616 & 0.639 \\
\bottomrule
\end{tabular}
\end{table}

%% -------------------------------------------------------------------------------------------
\section{ROC curves, calibration and clinical utility}

Among the frozen-feature image and fusion models, Brier scores ranged from 0.188 (gated fusion, DINOv2) to 0.226 (stacking, ConvNeXt); the phone-only model scored 0.182. Class-balanced training made image and fusion models overestimate risk, and decision-curve analysis \citep{vickers2006dca} showed no net benefit over treating everyone at thresholds up to 0.20; at 0.30 the stacked model gave a net benefit of 0.024 versus $-0.030$ for treat-all, but the phone-only model gave 0.082 (Fig.~\ref{fig:s-calib}).

\begin{figure}[!htbp]
\centering
\figfile{figures/fig2_roc.pdf}{0.55\textwidth}
\caption{Mean ROC curves ($\pm$1 SD over five repeats) for stacked fusion (DINOv2), the frozen DINOv2 image model, the phone-only model and the LightGBM symptom model. Dashed line: chance.}
\label{fig:s-roc}
\end{figure}

\begin{figure}[!htbp]
\centering
\figfile{figures/fig6a_calibration.pdf}{0.45\textwidth}\hfill
\figfile{figures/fig6b_decision_curve.pdf}{0.52\textwidth}
\caption{(A) Calibration curves and (B) decision curves for the main models. Treat-all and treat-none are shown for reference.}
\label{fig:s-calib}
\end{figure}

%% -------------------------------------------------------------------------------------------
\section{Additional shortcut controls and label dependence}

\paragraph{Large images} Restricting to photographs with an original short side of at least 500 pixels did not change within-phone AUC (DINOv2: 0.568 Samsung, 0.611 Xiaomi), so very small crops do not explain the shortcut.

\paragraph{Label dependence} Models discriminated better on unanimous than on split votes (DINOv2 image: 0.719 vs 0.678; symptoms: 0.623 vs 0.580). Image predictions correlated more with the physicians' vote share than symptom predictions did (Spearman 0.32 vs 0.13), though the phone-only model also reached 0.28. In a residual test, DINOv2 embeddings predicted the out-of-fold errors of the symptom model (ridge regression, 5-fold cross-validated $r = 0.29$; permutation $p = 0.002$ with 500 permutations), but this test cannot separate throat appearance from phone identity.

%% -------------------------------------------------------------------------------------------
\section{Explainability and subgroups}

We computed SHAP values \citep{lundberg2017shap} for a LightGBM symptom model and Grad-CAM maps \citep{selvaraju2017gradcam} from a ConvNeXt-Tiny fine-tuned on all patients for the median best epoch count (seven), used only for visualisation. Cough, age, vomiting, myalgia and sore throat had the largest mean absolute SHAP values (Fig.~\ref{fig:s-shap}). Grad-CAM maps of the most confident bacterial cases highlighted white spots in close-up photographs, whereas the most confident non-bacterial cases were wide views of the oral cavity with the uvula and tongue visible (Fig.~\ref{fig:s-gradcam}); framing and zoom may therefore act as an additional site cue. The image model discriminated better in children aged 14 or younger (AUC 0.77) than in adults aged 45 or older (0.64), and similarly in women (0.68) and men (0.72) (Fig.~\ref{fig:s-subgroups}).

\begin{figure}[!htbp]
\centering
\figfile{figures/figS2_shap.pdf}{0.75\textwidth}
\caption{SHAP summary of the LightGBM symptom model fitted on all patients.}
\label{fig:s-shap}
\end{figure}

\begin{figure}[!htbp]
\centering
\figfile{figures/fig5_gradcam.png}{\textwidth}
\caption{Grad-CAM maps of a ConvNeXt-Tiny fine-tuned on all patients for the six most confident bacterial (top) and non-bacterial (bottom) cases. Red marks high attention. Bacterial predictions are close-ups; non-bacterial predictions are wide views.}
\label{fig:s-gradcam}
\end{figure}

\begin{figure}[!htbp]
\centering
\figfile{figures/figS1_subgroups.pdf}{0.75\textwidth}
\caption{AUC of the frozen DINOv2 image model by sex, age band and phone, with 95\% bootstrap confidence intervals. Dashed line: overall AUC.}
\label{fig:s-subgroups}
\end{figure}

\bibliographystyle{unsrtnat}
\bibliography{references}

\end{document}
